% Nallino, Part II p. 224: Ptolemy's lunar model, redrawn from the print. The coordinates are those of a 400 dpi
% render of the page (y downward); E is the centre of the oblique circle (the Earth), F the centre of the
% excentric, F' its reflection in E, A and P the apogee and the perigee of the excentric, M a third position of
% the centre of the epicycle. The segment of the line of apsides inside the epicycle at A is dashed in the print.
\begin{tikzpicture}[x=0.07mm,y=-0.07mm,line width=0.4pt,every node/.style={font=\bfseries\itshape\small,inner sep=1pt}]
\coordinate (E) at (955,678);
\coordinate (F) at (955,585);
\coordinate (Fp) at (955,771);
\coordinate (A) at (955,155);
\coordinate (P) at (955,1015);
\coordinate (M) at (656,894);
\coordinate (L) at (999,646);
\coordinate (Hup) at (847,177);
\coordinate (Hlow) at (855,970);
\draw (E) circle[radius=525];
\draw (F) circle[radius=430];
\draw (A) circle[radius=110];
\draw (P) circle[radius=110];
\draw (M) circle[radius=110];
\draw[dashed] (A) -- (955,265);
\draw (955,265) -- (955,1203);
\draw (811,10) -- (E);
\draw (Hup) -- (A);
\draw (F) -- (M);
\draw (L) -- (567,958);
\draw (F) -- (L);
\draw (Fp) -- (554,936);
\draw (E) -- (639,1097);
\draw (E) -- (827,1052);
\draw (Hlow) -- (P);
\node at (957,122) {A};
\node at (822,210) {H};
\node at (988,296) {G};
\node at (410,580) {B};
\node at (985,560) {F};
\node at (1030,632) {L};
\node at (1505,632) {D};
\node at (983,700) {E};
\node at (995,776) {F′};
\node at (780,812) {R};
\node at (660,855) {M};
\node at (994,870) {G};
\node at (525,925) {H};
\node at (558,985) {G};
\node at (752,1005) {K};
\node at (825,960) {H};
\node at (980,985) {P};
\node at (632,1140) {N};
\node at (975,1155) {S};
\node at (948,1238) {C};
\end{tikzpicture}
